\documentclass[11pt,a4paper]{article}
\usepackage{fontspec}
\usepackage{kotex}
\usepackage{geometry}
\usepackage{graphicx}
\usepackage{xcolor}
\usepackage{tcolorbox}
\usepackage{booktabs}
\usepackage{array}
\usepackage{longtable}
\usepackage{colortbl}
\usepackage{fancyhdr}
\usepackage{titlesec}
\usepackage{hyperref}
\usepackage{emoji}
\usepackage{amsmath}
\usepackage{amssymb}
\usepackage{enumitem}

\geometry{left=25mm, right=25mm, top=30mm, bottom=30mm}

\setmainfont{Noto Sans CJK KR}
\setsansfont{Noto Sans CJK KR}
\setmonofont{Noto Sans CJK KR}
\setemojifont{Noto Color Emoji}

\definecolor{primary}{HTML}{1976D2}
\definecolor{darkbrown}{HTML}{3E2723}
\definecolor{mediumbrown}{HTML}{5D4037}
\definecolor{lightbrown}{HTML}{8D6E63}
\definecolor{cream}{HTML}{FFF3E0}
\definecolor{lightgray}{HTML}{F5F5F5}
\definecolor{tableheader}{HTML}{E8E8E8}

\titleformat{\section}
  {\Large\bfseries\color{primary}}
  {\thesection.}{0.5em}{}
\titleformat{\subsection}
  {\large\bfseries\color{darkbrown}}
  {\thesubsection}{0.5em}{}
\titleformat{\subsubsection}
  {\normalsize\bfseries\color{mediumbrown}}
  {\thesubsubsection}{0.5em}{}

\renewcommand{\contentsname}{목차}
\renewcommand{\figurename}{그림}
\renewcommand{\tablename}{표}

\tcbset{
  tipbox/.style={colback=cream,colframe=primary,boxrule=1.5pt,arc=3mm,
    left=5mm,right=5mm,top=3mm,bottom=3mm,fonttitle=\bfseries\color{primary}},
  formulabox/.style={colback=lightgray,colframe=lightgray,boxrule=0pt,arc=2mm,
    left=5mm,right=5mm,top=3mm,bottom=3mm},
  summarybox/.style={colback=white,colframe=primary,boxrule=2pt,arc=4mm,
    left=5mm,right=5mm,top=5mm,bottom=5mm}
}

\hypersetup{colorlinks=true,linkcolor=primary,urlcolor=primary,bookmarks=true,unicode=true}


\pagestyle{fancy}
\fancyhf{}
\fancyhead[R]{\small\color{lightbrown}마이크로비트 데이터로거 시스템 퀵가이드 v1.1.0.0}
\fancyfoot[C]{\thepage}
\renewcommand{\headrulewidth}{0.4pt}

\begin{document}

\begin{titlepage}
\centering
\vspace*{3cm}
{\Huge\bfseries\color{primary}\emoji{memo} 마이크로비트 데이터로거 시스템}\\[0.5cm]
{\Huge\bfseries\color{primary} 퀵가이드}
\\[1.5cm]
{\Large 시뮬레이터 버전: v1.0.3.0}
\\[2cm]
\begin{tcolorbox}[summarybox]
\centering
{\large 이 퀵가이드는 마이크로비트 데이터로거 시스템와 함께 보는 내비게이션입니다.}
\end{tcolorbox}
\vfill
{\large 사물인터넷과 빅데이터}\\[0.3cm]
{\large 청주교육대학교 컴퓨터교육과}
\vspace{1cm}
\end{titlepage}

\section*{\emoji{open-book} 이 문서 사용법}

\begin{tcolorbox}[tipbox]
\textbf{이 퀵가이드는 데이터로거 다이어그램과 함께 보는 내비게이션입니다.}

이 자료는 \textbf{정적 인포그래픽}입니다 — 클릭하거나 조작하는 기능은 없고, 보면서 이해하는 참고 자료입니다.

\end{tcolorbox}

\begin{center}
\renewcommand{\arraystretch}{1.4}
\begin{tabular}{|p{8.4cm}|p{6.6cm}|}
\hline
\rowcolor{tableheader}
\textbf{상황} & \textbf{참조 섹션} \\
\hline
전체 흐름 파악 & \emoji{high-voltage} 초속성 실습 \\
\hline
화면 배치가 궁금함 & \emoji{desktop-computer} 화면 구성 \\
\hline
센서 종류 확인 & \emoji{magnifying-glass-tilted-left} 관찰 포인트 \\
\hline
글자가 깨져요 & \emoji{sos} 트러블슈팅 \\
\hline
용어 모름 & \emoji{key} 핵심 용어 \\
\hline
\end{tabular}
\end{center}


\section*{\emoji{high-voltage} 1분 초속성 실습}

\begin{center}
\renewcommand{\arraystretch}{1.4}
\begin{tabular}{|p{1.5cm}|p{8.2cm}|p{5.3cm}|}
\hline
\rowcolor{tableheader}
\textbf{단계} & \textbf{행동} & \textbf{확인} \\
\hline
1 & 화면 왼쪽의 \textbf{마이크로비트 보드} 다이어그램을 찾으세요 & \emoji{green-circle} 초록 테두리 센서 카드 6개 보이나요? \\
\hline
2 & 오른쪽의 \textbf{Grove 확장 센서} 영역을 찾으세요 & \emoji{blue-circle} 파랑 테두리 센서 카드 10개 보이나요? \\
\hline
3 & 가운데 \textbf{데이터 저장 경로}(SD카드 \textrightarrow{} CSV \textrightarrow{} 엑셀) 흐름을 따라가세요 & 화살표가 보이면 \emoji{party-popper} \\
\hline
\end{tabular}
\end{center}


\section*{\emoji{rocket} 빠른 시작 — 이렇게 읽으세요}

\subsection*{\emoji{keycap-1} 데이터 수집 흐름 (위\textrightarrow{}아래)}

\begin{tcolorbox}[formulabox]
\begin{verbatim}
센서 → 마이크로비트 → 데이터로거(RTC+SD) → CSV 파일 → PC에서 분석
\end{verbatim}
\end{tcolorbox}

이 다이어그램은 이 흐름의 \textbf{첫 단계(센서 종류)}를 보여줍니다.


\subsection*{\emoji{keycap-2} 센서 2종류 구분하기}

\begin{center}
\renewcommand{\arraystretch}{1.4}
\begin{tabular}{|p{2.6cm}|p{2.4cm}|p{4.8cm}|p{5.2cm}|}
\hline
\rowcolor{tableheader}
\textbf{구분} & \textbf{색상} & \textbf{위치} & \textbf{특징} \\
\hline
\emoji{green-circle} 내장 센서 & 초록 테두리 & 좌측 (마이크로비트 안) & 추가 장비 없이 바로 사용 \\
\hline
\emoji{blue-circle} 외장 센서 & 파랑 테두리 & 우측 (Grove Shield) & 더 정밀·다양한 측정 \\
\hline
\emoji{purple-circle} 고정 장치 & 보라 테두리 & — & OLED 디스플레이, RTC 등 \\
\hline
\end{tabular}
\end{center}

\subsection*{\emoji{keycap-3} 데이터 항목 읽기}

각 센서 카드 아래의 \textbf{태그}가 측정 가능한 데이터입니다:


\begin{center}
\renewcommand{\arraystretch}{1.4}
\begin{tabular}{|p{3.4cm}|p{5.0cm}|p{6.6cm}|}
\hline
\rowcolor{tableheader}
\textbf{센서 예시} & \textbf{데이터 태그} & \textbf{의미} \\
\hline
SHT40 & \texttt{정밀 온도} \texttt{습도} & 이 센서로 온도와 습도를 측정 가능 \\
\hline
BMP280 & \texttt{기압} \texttt{고도} & 이 센서로 기압과 고도를 측정 가능 \\
\hline
마이크로비트 내장 & \texttt{온도} \texttt{조도} \texttt{가속도} 등 & 추가 장비 없이 측정 가능한 항목 \\
\hline
\end{tabular}
\end{center}


\section*{\emoji{desktop-computer} 화면 구성}

\begin{center}
\includegraphics[width=0.95\textwidth]{images/S3_DataLogger_screen.png}
\end{center}


\section*{\emoji{magnifying-glass-tilted-left} 관찰 포인트}

\begin{center}
\renewcommand{\arraystretch}{1.4}
\begin{tabular}{|p{2.8cm}|p{4.9cm}|p{7.3cm}|}
\hline
\rowcolor{tableheader}
\textbf{관찰 대상} & \textbf{찾아볼 장면} & \textbf{의미} \\
\hline
내장 센서 6개 & 좌측 초록 카드 & "마이크로비트만으로 이것들을 측정할 수 있다" \\
\hline
외장 센서 10개 & 우측 파랑 카드 & "Grove로 확장하면 더 많은 것을 측정" \\
\hline
저장 경로 & SD카드 \textrightarrow{} CSV \textrightarrow{} 엑셀 & "데이터는 파일로 저장되어 나중에 분석" \\
\hline
데이터 요약 수치 & 16종 센서, 24+ 데이터, 848 조합 & "센서를 조합하면 엄청 다양한 데이터 수집 가능" \\
\hline
센서별 데이터 태그 & 카드 아래 배지들 & "같은 센서라도 여러 종류의 데이터를 측정" \\
\hline
\end{tabular}
\end{center}


\section*{\emoji{sos} 도와주세요!}

\begin{center}
\renewcommand{\arraystretch}{1.4}
\begin{tabular}{|p{4.5cm}|p{3.9cm}|p{6.6cm}|}
\hline
\rowcolor{tableheader}
\textbf{증상} & \textbf{원인} & \textbf{해결책} \\
\hline
일부 한글이 깨져 보여요 & 브라우저 인코딩 설정 & Chrome에서 열기 + UTF-8 확인 \\
\hline
글자가 너무 작아요 & 화면 해상도 & \texttt{Ctrl + +}로 확대 (최대 150\% 추천) \\
\hline
오른쪽 센서가 잘려요 & 화면 너비 부족 & 브라우저를 전체 화면(F11)으로 전환 \\
\hline
인쇄가 잘 안 돼요 & 다크 테마 배경색 & 인쇄 설정에서 "배경 그래픽" 체크 \\
\hline
\end{tabular}
\end{center}


\section*{\emoji{key} 핵심 용어}

\begin{center}
\renewcommand{\arraystretch}{1.4}
\begin{tabular}{|p{1.5cm}|p{3.4cm}|p{4.5cm}|p{5.6cm}|}
\hline
\rowcolor{tableheader}
\textbf{용어} & \textbf{학술 용어} & \textbf{의미} & \textbf{쉬운 설명} \\
\hline
데이터로거 & Data Logger & 데이터 기록 장치 & 센서 값을 시간표와 함께 파일로 저장 \\
\hline
RTC & Real-Time Clock & 실시간 시계 모듈 & "몇 시 몇 분에 측정했는지" 기록 \\
\hline
SD카드 & SD Card & 저장 매체 & 데이터를 파일로 보관하는 메모리 카드 \\
\hline
CSV & Comma-Separated Values & 쉼표 구분 파일 & 엑셀에서 열 수 있는 데이터 형식 \\
\hline
내장 센서 & Built-in Sensor & 마이크로비트 기본 탑재 센서 & 추가 장비 없이 바로 측정 가능 \\
\hline
Grove 센서 & Grove Sensor & I2C로 추가 연결하는 외장 센서 & 더 정밀·다양한 측정이 가능 \\
\hline
\end{tabular}
\end{center}


\section*{\emoji{thought-balloon} 생각해보기}

\begin{tcolorbox}[summarybox]
\textbf{센서가 측정한 데이터는 어디에 저장되나요?}

\vspace{0.3em}
마이크로비트 메모리? SD카드? 클라우드?

데이터를 저장하지 않으면 어떤 문제가 생길까요?

\vspace{0.3em}
\textbf{프로젝트에 꼭 외장 센서가 필요할까요?}

\vspace{0.3em}
내장 센서 6종만으로도 온도·조도·가속도를 측정할 수 있습니다.

"내장 센서로 충분한가?"를 먼저 판단하면 불필요한 비용을 줄일 수 있어요.

\end{tcolorbox}


\section*{참고자료}

\begin{center}
\renewcommand{\arraystretch}{1.4}
\begin{tabular}{|p{1.5cm}|p{7.6cm}|p{2.7cm}|p{3.2cm}|}
\hline
\rowcolor{tableheader}
\textbf{\#} & \textbf{자료명} & \textbf{유형} & \textbf{비고} \\
\hline
1 & 마이크로비트데이터로거\_v1.0.3.0 & HTML 파일 & 다이어그램 \\
\hline
2 & 마이크로비트데이터로거\_개발참조\_v1.0.0.0 & 문서 & 기술 상세 \\
\hline
3 & 코딩관리지침\_v4.0.0.0 \S{}17 & 프로젝트 문서 & 퀵가이드 작성법 \\
\hline
\end{tabular}
\end{center}


\vfill
\begin{center}
\small\color{lightbrown}
\textit{마이크로비트 데이터로거 시스템 퀵가이드}
\end{center}
\end{document}